\documentclass[conference]{IEEEtran}
\IEEEoverridecommandlockouts
\usepackage{cite}
\usepackage{amsmath,amssymb,amsfonts}
\usepackage{algorithmic}
\usepackage{graphicx}
\usepackage{textcomp}
\usepackage{xcolor}
\usepackage{svg}
\usepackage{booktabs}
\usepackage{multirow}

\usepackage[utf8]{inputenc}
\usepackage[spanish]{babel}
\addto\captionsspanish{\renewcommand{\tablename}{Tabla}}

\usepackage{listings}
\usepackage{xcolor}

\renewcommand{\lstlistingname}{Código}

\lstset{
  language=Python,
  basicstyle=\ttfamily\small,
  frame=single,
  breaklines=true,
  keywordstyle=\color{blue},
  commentstyle=\color{gray},
  stringstyle=\color{teal},
  numbers=left,
  numberstyle=\tiny,
  captionpos=b
}

\usepackage{caption}
\captionsetup[table]{skip=10pt}

\title{Laboratorio 03: Clasificadores fundamentales \\
\vspace{0.5em} \huge \textbf{Asignatura}: \textit{Machine Learning} (2026-01)\\
\vspace{0.5em} \large \textbf{Profesor}: Dr. Juan Bekios Calfa}

\author{
\IEEEauthorblockN{1\textsuperscript{st} Nombre alumno 1}
\IEEEauthorblockA{\textit{Escuela de Ingeniería} \\
\textit{Universidad Católica del Norte}\\
Carrera\\
\texttt{correo01@electronico.ucn.cl}}
\and
\IEEEauthorblockN{2\textsuperscript{nd} Nombre alumno 2}
\IEEEauthorblockA{\textit{Escuela de Ingeniería} \\
\textit{Universidad Católica del Norte}\\
Carrera\\
\texttt{correo02@electronico.ucn.cl}}
\and
\IEEEauthorblockN{3\textsuperscript{rd} Nombre alumno 3}
\IEEEauthorblockA{\textit{Escuela de Ingeniería} \\
\textit{Universidad Católica del Norte}\\
Carrera\\
\texttt{correo03@electronico.ucn.cl}}
}

\begin{document}

\maketitle

\begin{abstract}
\noindent
Este informe presenta los resultados del laboratorio de clasificadores fundamentales aplicados a la predicción de niveles de deterioro cognitivo a partir de 15 atributos binarios. Se implementaron cinco clasificadores ---Regresión Logística, SVM lineal, SVM con kernel RBF, Árbol de decisión y K-NN--- sobre seis variables objetivo derivadas de la escala GDS. Se utilizó validación cruzada anidada con ajuste de hiperparámetros mediante GridSearchCV para evitar sesgo de selección. Los resultados muestran que SVM con kernel RBF obtiene el mejor desempeño en la mayoría de los objetivos, aunque K-NN y Regresión Logística son competitivos en problemas con menor número de clases.
\end{abstract}

\section{Introducción}

El deterioro cognitivo es un problema de creciente relevancia en poblaciones adultas mayores. Su detección temprana permite implementar intervenciones oportunas que pueden mejorar la calidad de vida de los pacientes. La escala GDS (\textit{Global Deterioration Scale}) es un instrumento ampliamente utilizado para clasificar el nivel de deterioro cognitivo en siete etapas.

En este laboratorio se trabaja con un conjunto de 1119 observaciones, cada una descrita por 15 atributos binarios que representan respuestas a un test neuropsicológico. El objetivo es construir, ajustar y comparar clasificadores fundamentales de Machine Learning para predecir distintas agrupaciones de la escala GDS.

Se definen seis experimentos independientes, uno por cada columna objetivo: \texttt{GDS}, \texttt{GDS\_R1}, \texttt{GDS\_R2}, \texttt{GDS\_R3}, \texttt{GDS\_R4} y \texttt{GDS\_R5}. Cada objetivo representa una forma distinta de agrupar los niveles de deterioro cognitivo, con distinto número de clases y distinto grado de desbalance.

Los clasificadores implementados son: Regresión Logística, SVM lineal, SVM con kernel RBF, Árbol de decisión y K-NN. Todos se implementan usando \texttt{Pipeline} y \texttt{GridSearchCV} de scikit-learn, siguiendo un protocolo de validación cruzada anidada que separa la selección de hiperparámetros de la evaluación final.

El informe se organiza en las siguientes secciones: marco teórico con los fundamentos de cada clasificador, metodología detallando el diseño experimental, resultados con tablas comparativas, discusión de los hallazgos y conclusiones del trabajo.

\section{Marco Teórico}

\subsection{Regresión Logística}
La regresión logística es un modelo lineal para clasificación probabilística. Estima la probabilidad de pertenencia a cada clase mediante la función softmax:
\begin{equation}
P(y = k \mid \mathbf{x}) = \frac{e^{\mathbf{w}_k^T \mathbf{x} + b_k}}{\sum_{j} e^{\mathbf{w}_j^T \mathbf{x} + b_j}}
\end{equation}
El hiperparámetro \(C\) controla la regularización inversa: valores pequeños de \(C\) aplican mayor regularización, lo que ayuda a prevenir el sobreajuste en problemas con muchas variables.

\subsection{SVM Lineal}
\textit{Support Vector Machine} con kernel lineal busca el hiperplano que maximiza el margen entre clases. La formulación primal es:
\begin{equation}
\min_{\mathbf{w}, b} \frac{1}{2} \|\mathbf{w}\|^2 + C \sum_{i=1}^{n} \xi_i
\end{equation}
El parámetro \(C\) controla el equilibrio entre margen amplio y penalización por errores de clasificación.

\subsection{SVM con Kernel RBF}
El kernel RBF (\textit{Radial Basis Function}) permite construir fronteras no lineales midiendo similitud local entre observaciones:
\begin{equation}
K(\mathbf{x}_i, \mathbf{x}_j) = \exp(-\gamma \|\mathbf{x}_i - \mathbf{x}_j\|^2)
\end{equation}
El parámetro \(\gamma\) controla la influencia local de cada punto: valores altos generan fronteras más complejas.

\subsection{Árbol de Decisión}
El árbol de decisión divide el espacio de atributos mediante reglas binarias del tipo \texttt{atributo <= umbral}. Es altamente interpretable, especialmente con variables binarias. Los hiperparámetros de regularización incluyen \texttt{max\_depth} (profundidad máxima) y \texttt{min\_samples\_leaf} (mínimo de muestras por hoja).

\subsection{K-NN}
K-NN clasifica una nueva observación según la clase mayoritaria entre sus \(k\) vecinos más cercanos. Es sensible al escalamiento de los datos y a la elección de \(k\), la métrica de distancia (euclidiana o Manhattan) y el tipo de ponderación (uniforme o por distancia).

\section{Metodología}

\subsection{Dataset y preprocesamiento}
El dataset contiene 1119 observaciones con 15 atributos binarios de entrada: \texttt{Día, Mes, Año, Estación, País, Ciudad, CalleLugar, NumeroPiso, Miguel2, González2, Avenida2, Imperial2, A682, Caldera2, Copiapo2}. Se excluye el identificador \texttt{ID} como predictor. Las columnas objetivo son \texttt{GDS, GDS\_R1, GDS\_R2, GDS\_R3, GDS\_R4, GDS\_R5}. Los atributos se convierten a tipo \texttt{float} y las etiquetas a \texttt{int}.

\subsection{Validación cruzada anidada}
Para evitar el sesgo de selección de hiperparámetros, se implementa validación cruzada anidada:

\begin{itemize}
  \item \textbf{Ciclo externo:} Estima el desempeño final del modelo usando \texttt{StratifiedKFold} con \(k_{\text{externo}} = \min(5, n_{\text{mín}})\), donde \(n_{\text{mín}}\) es el soporte de la clase menos frecuente.
  \item \textbf{Ciclo interno:} Selecciona los hiperparámetros óptimos mediante \texttt{GridSearchCV} con scoring \texttt{f1\_macro}. Se usa \(k_{\text{interno}} = \max(2, \min(3, k_{\text{externo}}))\).
\end{itemize}

Cuando una clase queda con un solo ejemplo en el entrenamiento externo, el ciclo interno usa \texttt{KFold} no estratificado y se registra una advertencia metodológica.

\subsection{Pipelines y grillas de hiperparámetros}
Cada clasificador se implementa dentro de un \texttt{Pipeline} de scikit-learn. Los modelos con escalamiento incluyen \texttt{StandardScaler}. La Tabla~\ref{tab:grids} resume las grillas exploradas.

\begin{table}[h]
\centering
\caption{Grillas de hiperparámetros por modelo}
\label{tab:grids}
\begin{tabular}{lll}
\toprule
\textbf{Modelo} & \textbf{Parámetro} & \textbf{Valores} \\
\midrule
Regresión Logística & \(C\) & 0.01, 0.1, 1, 10 \\
SVM lineal & \(C\) & 0.01, 0.1, 1, 10 \\
SVM RBF & \(C\) & 0.1, 1, 10 \\
 & \(\gamma\) & scale, 0.01, 0.1, 1 \\
Árbol de decisión & max\_depth & 2, 3, 4, 5, None \\
 & min\_samples\_leaf & 1, 3, 5, 10 \\
K-NN & n\_neighbors & 3, 5, 7, 9, 11 \\
 & weights & uniform, distance \\
 & metric & euclidean, manhattan \\
\bottomrule
\end{tabular}
\end{table}

\subsection{Métricas de evaluación}
Se reportan las siguientes métricas para cada modelo: F1 macro (promedio no ponderado del F1 por clase), balanced accuracy, precisión macro, recall macro y estabilidad (\(1 - \text{std(F1 macro)}\)). Adicionalmente, se calcula el Índice Comparativo Normalizado (ICN) que combina las métricas mediante ponderación:
\begin{equation}
\text{ICN} = 0.40 \cdot \hat{F1}_{\text{macro}} + 0.25 \cdot \hat{BA} + 0.20 \cdot \hat{R}_{\text{macro}} + 0.10 \cdot \hat{P}_{\text{macro}} + 0.05 \cdot \hat{E}
\end{equation}
donde cada métrica se normaliza entre 0 y 1 mediante normalización min-max.

\section{Resultados}

\subsection{Distribución de clases}
La Tabla~\ref{tab:distribucion} muestra la distribución de clases para cada objetivo. Se observa un fuerte desbalance en todos los experimentos, siendo \texttt{GDS} el caso más extremo con 7 clases y una clase con solo 2 ejemplos.

\begin{table}[h]
\centering
\caption{Distribución de clases por objetivo}
\label{tab:distribucion}
\begin{tabular}{lcccc}
\toprule
\textbf{Objetivo} & \textbf{Clases} & \textbf{n$_{\text{mín}}$} & \textbf{k$_{\text{ext}}$} \\
\midrule
GDS & \{1:149, 2:500, 3:298, 4:108, 5:42, 6:20, 7:2\} & 2 & 2 \\
GDS\_R1 & \{1:947, 2:150, 3:22\} & 22 & 5 \\
GDS\_R2 & \{1:649, 2:298, 3:172\} & 172 & 5 \\
GDS\_R3 & \{1:947, 3:172\} & 172 & 5 \\
GDS\_R4 & \{1:149, 2:906, 3:64\} & 64 & 5 \\
GDS\_R5 & \{1:149, 2:798, 3:172\} & 149 & 5 \\
\bottomrule
\end{tabular}
\end{table}

\subsection{Resultados por experimento}
A continuación se presentan las tablas de resultados para cada uno de los seis experimentos. Se reporta el promedio y desviación estándar del F1 macro en los folds externos, junto con las métricas complementarias y el ICN.

\subsubsection*{Experimento GDS (7 clases)}
La Tabla~\ref{tab:gds} muestra los resultados para \texttt{GDS}. Con \(k=2\) folds externos debido a \(n_{\text{mín}}=2\), las estimaciones tienen alta variabilidad.

\begin{table}[h]
\centering
\caption{Resultados para GDS}
\label{tab:gds}
\begin{tabular}{lccccc}
\toprule
\textbf{Modelo} & \textbf{F1 macro} & \textbf{Bal. Acc.} & \textbf{ICN} \\
\midrule
Regresión Logística & 0.295 $\pm$ 0.038 & 0.440 & 0.660 \\
SVM lineal & 0.308 $\pm$ 0.023 & 0.470 & 0.843 \\
SVM RBF & \textbf{0.335 $\pm$ 0.039} & 0.452 & \textbf{0.872} \\
Árbol decisión & 0.223 $\pm$ 0.033 & 0.412 & 0.263 \\
K-NN & 0.329 $\pm$ 0.014 & 0.338 & 0.529 \\
\bottomrule
\end{tabular}
\end{table}

\subsubsection*{Experimento GDS\_R1 (3 clases)}
La Tabla~\ref{tab:gds_r1} muestra los resultados para \texttt{GDS\_R1}. Con 3 clases y \(k=5\) folds, las estimaciones son más estables.

\begin{table}[h]
\centering
\caption{Resultados para GDS\_R1}
\label{tab:gds_r1}
\begin{tabular}{lccccc}
\toprule
\textbf{Modelo} & \textbf{F1 macro} & \textbf{Bal. Acc.} & \textbf{ICN} \\
\midrule
Regresión Logística & 0.652 $\pm$ 0.062 & 0.763 & 0.644 \\
SVM lineal & 0.654 $\pm$ 0.056 & 0.770 & 0.708 \\
SVM RBF & \textbf{0.671 $\pm$ 0.058} & \textbf{0.781} & \textbf{0.827} \\
Árbol decisión & 0.592 $\pm$ 0.051 & 0.708 & 0.182 \\
K-NN & 0.690 $\pm$ 0.065 & 0.678 & 0.500 \\
\bottomrule
\end{tabular}
\end{table}

\subsubsection*{Experimento GDS\_R2 (3 clases)}
La Tabla~\ref{tab:gds_r2} muestra los resultados para \texttt{GDS\_R2}. Las tres clases tienen soporte suficiente para una validación estable.

\begin{table}[h]
\centering
\caption{Resultados para GDS\_R2}
\label{tab:gds_r2}
\begin{tabular}{lccccc}
\toprule
\textbf{Modelo} & \textbf{F1 macro} & \textbf{Bal. Acc.} & \textbf{ICN} \\
\midrule
Regresión Logística & 0.663 $\pm$ 0.043 & 0.666 & 0.769 \\
SVM lineal & 0.667 $\pm$ 0.048 & 0.662 & 0.804 \\
SVM RBF & \textbf{0.675 $\pm$ 0.056} & \textbf{0.667} & \textbf{0.950} \\
Árbol decisión & 0.648 $\pm$ 0.032 & 0.658 & 0.512 \\
K-NN & 0.639 $\pm$ 0.031 & 0.622 & 0.099 \\
\bottomrule
\end{tabular}
\end{table}

\subsubsection*{Experimento GDS\_R3 (binario)}
La Tabla~\ref{tab:gds_r3} muestra los resultados para \texttt{GDS\_R3}. Al ser un problema binario con soporte suficiente, todos los modelos alcanzan buen desempeño.

\begin{table}[h]
\centering
\caption{Resultados para GDS\_R3}
\label{tab:gds_r3}
\begin{tabular}{lccccc}
\toprule
\textbf{Modelo} & \textbf{F1 macro} & \textbf{Bal. Acc.} & \textbf{ICN} \\
\midrule
Regresión Logística & 0.764 $\pm$ 0.031 & 0.842 & 0.714 \\
SVM lineal & 0.757 $\pm$ 0.043 & 0.836 & 0.618 \\
SVM RBF & \textbf{0.792 $\pm$ 0.039} & \textbf{0.838} & \textbf{0.907} \\
Árbol decisión & 0.733 $\pm$ 0.077 & 0.795 & 0.218 \\
K-NN & 0.783 $\pm$ 0.046 & 0.753 & 0.474 \\
\bottomrule
\end{tabular}
\end{table}

\subsubsection*{Experimento GDS\_R4 (3 clases)}
La Tabla~\ref{tab:gds_r4} muestra los resultados para \texttt{GDS\_R4}. Con desbalance moderado, K-NN obtiene el mejor ICN.

\begin{table}[h]
\centering
\caption{Resultados para GDS\_R4}
\label{tab:gds_r4}
\begin{tabular}{lccccc}
\toprule
\textbf{Modelo} & \textbf{F1 macro} & \textbf{Bal. Acc.} & \textbf{ICN} \\
\midrule
Regresión Logística & 0.532 $\pm$ 0.015 & \textbf{0.718} & \textbf{0.596} \\
SVM lineal & 0.521 $\pm$ 0.025 & 0.709 & 0.494 \\
SVM RBF & \textbf{0.550 $\pm$ 0.039} & 0.590 & 0.248 \\
Árbol decisión & 0.515 $\pm$ 0.027 & 0.659 & 0.273 \\
K-NN & 0.592 $\pm$ 0.040 & 0.588 & 0.500 \\
\bottomrule
\end{tabular}
\end{table}

\subsubsection*{Experimento GDS\_R5 (3 clases)}
La Tabla~\ref{tab:gds_r5} muestra los resultados para \texttt{GDS\_R5}. SVM RBF lidera tanto en F1 macro como en ICN.

\begin{table}[h]
\centering
\caption{Resultados para GDS\_R5}
\label{tab:gds_r5}
\begin{tabular}{lccccc}
\toprule
\textbf{Modelo} & \textbf{F1 macro} & \textbf{Bal. Acc.} & \textbf{ICN} \\
\midrule
Regresión Logística & 0.521 $\pm$ 0.029 & 0.657 & 0.527 \\
SVM lineal & 0.507 $\pm$ 0.030 & 0.655 & 0.431 \\
SVM RBF & \textbf{0.549 $\pm$ 0.011} & \textbf{0.664} & \textbf{0.739} \\
Árbol decisión & 0.510 $\pm$ 0.046 & 0.621 & 0.284 \\
K-NN & 0.584 $\pm$ 0.049 & 0.571 & 0.500 \\
\bottomrule
\end{tabular}
\end{table}

\subsection{Resumen global}
La Tabla~\ref{tab:resumen} presenta el mejor modelo por objetivo según ICN.

\begin{table}[h]
\centering
\caption{Mejor modelo por objetivo}
\label{tab:resumen}
\begin{tabular}{lccc}
\toprule
\textbf{Objetivo} & \textbf{Mejor modelo} & \textbf{F1 macro} & \textbf{ICN} \\
\midrule
GDS & SVM RBF & 0.335 & 0.872 \\
GDS\_R1 & SVM RBF & 0.671 & 0.827 \\
GDS\_R2 & SVM RBF & 0.675 & 0.950 \\
GDS\_R3 & SVM RBF & 0.792 & 0.907 \\
GDS\_R4 & Regresión Logística & 0.532 & 0.596 \\
GDS\_R5 & SVM RBF & 0.549 & 0.739 \\
\bottomrule
\end{tabular}
\end{table}

\section{Discusión}

Los resultados muestran que SVM con kernel RBF es el clasificador con mejor desempeño general, ganando en 5 de los 6 experimentos según el ICN. Esto sugiere que las fronteras de decisión entre las clases de deterioro cognitivo no son linealmente separables, y el kernel RBF captura relaciones no lineales que benefician la clasificación.

\noindent \textbf{Caso GDS:} Es el experimento más desafiante con 7 clases y una clase con solo 2 ejemplos. Todos los modelos obtienen F1 macro por debajo de 0.34, lo que refleja la dificultad de clasificar con soporte tan bajo. Las advertencias metodológicas indican que el ciclo interno usó KFold no estratificado debido a clases con un solo ejemplo en entrenamiento.

\noindent \textbf{Caso GDS\_R4:} Es el único donde la Regresión Logística supera a SVM RBF en ICN. Esto puede deberse a que la clase mayoritaria (clase 2 con 906 ejemplos) domina el problema, y un modelo lineal con regularización fuerte (\(C = 0.01\)) es suficiente para obtener buen desempeño.

\noindent \textbf{Estabilidad:} El Árbol de decisión muestra la mayor variabilidad entre folds (desviación estándar alta), especialmente en GDS\_R3 (std = 0.077), lo que indica sensibilidad a la composición del fold. En contraste, K-NN y Regresión Logística tienden a ser más estables.

\noindent \textbf{Hiperparámetros seleccionados:} La regresión logística y SVM lineal favorecen consistentemente \(C = 0.01\), el valor más regularizado, lo que sugiere que agregar regularización beneficia el desempeño. SVM RBF selecciona frecuentemente \(\gamma = 0.01\) o \texttt{scale}, indicando que una influencia local suave es preferible.

\noindent \textbf{Limitaciones:} El desbalance extremo (clases con 2--22 ejemplos) limita la confiabilidad de las estimaciones para clases minoritarias. Los resultados deben interpretarse como evidencia experimental, no como conclusiones clínicas.

\section{Conclusiones}

Se implementaron y compararon cinco clasificadores fundamentales sobre seis variables objetivo derivadas de la escala GDS, utilizando validación cruzada anidada para garantizar una estimación honesta del desempeño.

Los principales hallazgos son:
\begin{itemize}
  \item SVM con kernel RBF es el modelo más robusto, obteniendo el mejor ICN en 5 de 6 experimentos.
  \item El desbalance de clases afecta significativamente el desempeño, especialmente en el objetivo GDS con 7 clases.
  \item La regularización (\(C\) pequeño) beneficia consistentemente a los modelos lineales.
  \item El Árbol de decisión es el modelo menos estable, con alta variabilidad entre folds.
  \item En problemas con clases muy desbalanceadas, ningún modelo supera F1 macro de 0.35.
\end{itemize}

Como trabajo futuro, se sugiere explorar técnicas de balanceo como SMOTE, evaluar la agrupación de clases minoritarias, y probar modelos más complejos como Gradient Boosting o Random Forest.

\section{Código y Reproducibilidad}

\subsection{Estructura del proyecto}
El código se organiza de forma modular en los siguientes archivos:

\begin{itemize}
  \item \texttt{src/main.py}: orquestación de experimentos y generación de reportes.
  \item \texttt{src/models.py}: definición de pipelines y grillas de hiperparámetros.
  \item \texttt{src/evaluation.py}: implementación de validación cruzada anidada.
  \item \texttt{src/data\_loader.py}: carga del archivo .sav y preparación de datos.
  \item \texttt{src/reports.py}: generación de tablas LaTeX, PDF, CSVs y JSON.
  \item \texttt{src/settings.py}: configuración de rutas y parámetros.
  \item \texttt{config/paths.yaml}: configuración de rutas de dataset y salidas.
\end{itemize}

\subsection{Ejecución}
Para reproducir los experimentos, se debe crear el entorno conda y ejecutar:
\begin{lstlisting}[caption={Comandos para reproducir experimentos}, label={lst:repro}]
conda env create -f environment.yml
conda activate lab03_ml_2026_01
python src/main.py
\end{lstlisting}

Se puede ejecutar un subconjunto de objetivos:
\begin{lstlisting}[caption={Ejecutar objetivos específicos}, label={lst:subset}]
python src/main.py --targets GDS_R2 GDS_R5
\end{lstlisting}

\subsection{Buenas prácticas}
El código aplica las siguientes buenas prácticas: uso de \texttt{Pipeline} para evitar fuga de información, separación estricta entre selección de hiperparámetros y evaluación mediante validación anidada, configuración externalizada en YAML, generación automatizada de reportes, y manejo explícito de advertencias metodológicas.

\subsection{Salidas generadas}
Los resultados se almacenan en \texttt{outputs/}:
\begin{itemize}
  \item \texttt{tables/resultados\_experimentos.tex} y \texttt{.pdf}: tablas por experimento.
  \item \texttt{tables/resumen\_resultados.csv}: resumen tabular de todos los modelos.
  \item \texttt{tables/resultados\_detallados.json}: resultados completos en formato estructurado.
  \item \texttt{confusion\_matrices/}: matrices de confusión por modelo y objetivo.
  \item \texttt{per\_class/}: precisión, recall y F1 por clase.
  \item \texttt{advertencias.txt}: advertencias metodológicas de validación.
\end{itemize}

% Apéndices
\appendices

\section{Anexos}

\subsection{Hiperparámetros seleccionados por modelo}

\begin{table}[h]
\centering
\caption{Hiperparámetros más frecuentes por modelo y objetivo}
\label{tab:params}
\begin{tabular}{lll}
\toprule
\textbf{Objetivo} & \textbf{Modelo} & \textbf{Hiperparámetros} \\
\midrule
GDS & RL & C=0.01 (1/2) \\
 & SVM lineal & C=0.01 (1/2) \\
 & SVM RBF & C=10, $\gamma$=0.1 (1/2) \\
 & Árbol & max\_depth=5, min\_samples\_leaf=5 (1/2) \\
 & K-NN & metric=manhattan, n=7, weights=distance (1/2) \\
\midrule
GDS\_R1 & RL & C=0.01 (5/5) \\
 & SVM lineal & C=0.01 (3/5) \\
 & SVM RBF & C=0.1, $\gamma$=0.01 (4/5) \\
 & Árbol & max\_depth=3, min\_samples\_leaf=1 (2/5) \\
 & K-NN & metric=euclidean, n=11, weights=distance (2/5) \\
\midrule
GDS\_R2 & RL & C=0.01 (4/5) \\
 & SVM lineal & C=0.01 (4/5) \\
 & SVM RBF & C=0.1, $\gamma$=0.01 (4/5) \\
 & Árbol & max\_depth=4, min\_samples\_leaf=1 (1/5) \\
 & K-NN & metric=euclidean, n=7, weights=uniform (2/5) \\
\midrule
GDS\_R3 & RL & C=0.01 (5/5) \\
 & SVM lineal & C=0.01 (5/5) \\
 & SVM RBF & C=0.1, $\gamma$=0.01 (5/5) \\
 & Árbol & max\_depth=2, min\_samples\_leaf=1 (2/5) \\
 & K-NN & metric=euclidean, n=3, weights=uniform (2/5) \\
\midrule
GDS\_R4 & RL & C=0.01 (3/5) \\
 & SVM lineal & C=0.01 (4/5) \\
 & SVM RBF & C=0.1, $\gamma$=1 (3/5) \\
 & Árbol & max\_depth=None, min\_samples\_leaf=1 (4/5) \\
 & K-NN & metric=euclidean, n=3, weights=uniform (3/5) \\
\midrule
GDS\_R5 & RL & C=0.01 (4/5) \\
 & SVM lineal & C=0.01 (5/5) \\
 & SVM RBF & C=0.1, $\gamma$=scale (4/5) \\
 & Árbol & max\_depth=None, min\_samples\_leaf=1 (2/5) \\
 & K-NN & metric=euclidean, n=5, weights=uniform (2/5) \\
\bottomrule
\end{tabular}
\end{table}

\section{Rúbrica de Autoevaluación}

\begin{table}[h]
\centering
\begin{tabular}{|l|c|}
\hline
\textbf{Criterio} & \textbf{¿Cumplido?} \\
\hline
Introducción clara & Sí / No \\
Marco teórico adecuado & Sí / No \\
Metodología detallada & Sí / No \\
Resultados completos y visuales & Sí / No \\
Discusión crítica & Sí / No \\
Conclusiones reflexivas & Sí / No \\
Calidad y estructura del código & Sí / No \\
Presentación clara y sin errores & Sí / No \\
Reproducibilidad & Sí / No \\
Aportes adicionales & Sí / No \\
\hline
\end{tabular}
\caption{Tabla de autoevaluación}
\end{table}

\end{document}
